\subsection{Polar sets and orthogonal subspaces}

DEFINITION 2. --- Let F and G be two (real) vector spaces in duality. For every set M of F we call the polar of M, the set of those \(y \in G\) for which \(\langle x, y \rangle \geqslant -1\) for all \(x \in M\) . (For complex vector spaces, cf.~II, p.~64.)

If \(G_{1}\) , \(G_{2}\) are two subspaces of \(F^{*}\) such that \(G_{1} \subset G_{2}\) , then the polar of M in \(G_{1}\) is the intersection of \(G_{1}\) with the polar of M in \(G_{2}\) .

When there is no danger of confusion we use \(M^{\circ}\) to denote the polar, in G, of the subset M of F. Similarly we define the polar in F of a set in G.

Obviously, for every scalar \(\lambda \neq 0\) and all \(\mathbf{M} \subset \mathbf{F}\) , we have \((\lambda \mathbf{M})^{\circ} = \lambda^{-1} \mathbf{M}^{\circ}\) . The relation \(\mathbf{M} \subset \mathbf{N} \subset \mathbf{F}\) implies \(\mathbf{N}^{\circ} \subset \mathbf{M}^{\circ}\) ; if \(\mathbf{N}\) absorbs \(\mathbf{M}\) then \(\mathbf{M}^{\circ}\) absorbs \(\mathbf{N}^{\circ}\) ; for every family \((\mathbf{M}_{\alpha})\) of sets of \(\mathbf{F}\) , the polar set of \(\bigcup_{\alpha} \mathbf{M}_{\alpha}\) is the intersection of the polar sets \(\mathbf{M}_{\alpha}^{\circ}\) . Since, for \(y \in \mathbf{M}^{\circ}\) , the closed half-spaces defined by the relations \(\langle x, y \rangle \geqslant -1\) contain 0 and \(\mathbf{M}\) , we see that if \(\mathbf{M}_{1}\) is the convex envelope of \(\mathbf{M} \cup \{0\}\) , then \(\mathbf{M}_{1}^{\circ} = \mathbf{M}^{\circ}\) .

Clearly \(\mathbf{M} \subset \mathbf{M}^{\circ \circ}\) . Hence

\[
(\mathbf {M} ^ {\circ \circ}) ^ {\circ} \subset \mathbf {M} ^ {\circ} \subset (\mathbf {M} ^ {\circ}) ^ {\circ \circ} = (\mathbf {M} ^ {\circ \circ}) ^ {\circ}
\]

i.e.~\(\mathbf{M}^{\circ \circ \circ} = \mathbf{M}^{\circ}\) (cf.~S, III, § 1.5, prop. 2).

If \(\mathbf{M}\) is a symmetric subset of \(\mathrm{F}\) , \(\mathbf{M}^{\circ}\) is a symmetric subset of \(\mathbf{G}\) ; \(\mathbf{M}^{\circ}\) is also in this case the set of \(y \in \mathbf{G}\) such that \(|\langle x, y \rangle| \leqslant 1\) for all \(x \in \mathbf{M}\) .

PROPOSITION 4. --- (i) For any set M of F, the polar set \(\mathbf{M}^{\circ}\) is a convex set that contains 0 and is closed in G for the topology \(\sigma(\mathbf{G},\mathbf{F})\) .

\begin{enumerate}
\def\labelenumi{(\roman{enumi})}
\setcounter{enumi}{1}
\item
  If \(\mathbf{M}\) is a cone of vertex 0, then \(\mathbf{M}^{\circ}\) is a cone of vertex 0 and it is also the set of \(y\in \mathbf{G}\) such that \(\langle x,y\rangle \geqslant 0\) for all \(x\in \mathbf{M}\) .
\item
  If \(\mathbf{M}\) is a vector subspace of \(\mathbf{F}\) , then \(\mathbf{M}^{\circ}\) is a vector subspace of \(\mathbf{G}\) , and it is also the set of \(y \in \mathbf{G}\) such that \(\langle x, y \rangle = 0\) for all \(x \in \mathbf{M}\) .
\item
  Since the linear forms \(y \mapsto \langle x, y \rangle\) are continuous for \(\sigma(\mathbf{G}, \mathbf{F})\) the statement follows immediately from the definitions and the fact that a half-space determined by a hyperplane is convex.
\item
  If \(\mathbf{M}\) is a cone with vertex 0 and if \(x \in \mathbf{M}\) , \(y \in \mathbf{M}^{\circ}\) , then as \(\lambda x \in \mathbf{M}\) , for all \(\lambda > 0\) , we have \(\langle \lambda x, y \rangle \geqslant -1\) , i.e.~\(\lambda \langle x, y \rangle \geqslant -1\) . Since this holds for all \(\lambda > 0\) , it follows that \(\langle x, y \rangle \geqslant 0\) , and (ii) is proved.
\item
  Similarly, if \(\mathbf{M}\) is a vector subspace of \(\mathbf{F}\) , the relations \(x \in \mathbf{M}\) , \(y \in \mathbf{M}^{\circ}\) imply, this time, that \(\lambda \langle x, y \rangle \geqslant -1\) for all real \(\lambda\) which is possible only if \(\langle x, y \rangle = 0\) .
\end{enumerate}

If \(\mathbf{M}\) is a vector subspace of \(\mathbf{F}\) we say that \(\mathbf{M}^{\circ}\) is the orthogonal of \(\mathbf{M}\) in \(\mathbf{G}\) ; if \(\mathbf{G} \subset \mathbf{F}^{*}\) , then \(\mathbf{M}^{\circ}\) is the intersection of \(\mathbf{G}\) , and of the subspace orthogonal to \(\mathbf{M}\) in the algebraic dual \(\mathbf{F}^{*}\) of \(\mathbf{F}\) (A, II, § 2.4, def. 4).

For a vector subspace M of F and a vector subspace N of G we say that M and N are orthogonal if \(M \subset N^{0}\) (or, equivalently, if \(N \subset M^{\circ}\) ).

THEOREM 1 (The bipolar theorem). --- Let F, G be two real vector spaces in duality. For every subset M of F the polar set \(M^{\circ\circ}\) in F of the polar set \(M^{\circ}\) of M in G is the closed convex envelope (for \(\sigma(F, G)\) ) of \(M \cup \{0\}\) .

We have seen that we need only consider the case where M is convex and \(0 \in M\) . Denote the closure, in topology \(\sigma(F, G)\) of M by \(\overline{M}\) , then \(\overline{M}\) is a convex set in F; prop. 4 of II, p.~44 shows that \(M^{\circ\circ} \supset \overline{M}\) . On the other hand if \(a \in F\) does not belong to \(\overline{M}\) then there exists a closed hyperplane H in F which separates a strictly from \(\overline{M}\) (II, p.~38, prop. 4); since H does not contain 0, there exists \(y \in G\) such that H has the equation \(\langle x, y \rangle = -1\) (II, p.~43, prop. 3); thus \(\langle x, y \rangle > -1\) for all \(x \in \overline{M}\) and \(\langle a, y \rangle < -1\) . This implies that \(y \in M^{\circ}\) and \(a \notin M^{\circ\circ}\) , and the relation \(M^{\circ\circ} = \overline{M}\) follows.

COROLLARY 1. --- For any family \((\mathbf{M}_{\alpha})\) of closed convex sets of F (in the topology \(\sigma(\mathrm{F}, \mathrm{G})\) ), each containing 0, the polar set of the intersections \(M = \bigcap_{\alpha} M_{\alpha}\) is the convex closed envelope (for \(\sigma(\mathrm{G}, \mathrm{F})\) ) of the union of the \(M_{\alpha}^{\circ}\) .

For, if N is this convex closed envelope, then

\[
\mathbf {N} ^ {\circ} = \bigcap_ {\alpha} \mathbf {M} _ {\alpha} ^ {\circ \circ} = \bigcap_ {\alpha} \mathbf {M} _ {\alpha} = \mathbf {M}
\]

whence \(N = N^{\circ\circ} = M^{\circ}\) .

The conclusion of cor. 1 does not necessarily hold if the \(M_{\alpha}\) are not convex.

COROLLARY 2. --- For every vector subspace M of F, the subspace \(M^{\circ\circ}\) is the closure of M in the topology \(\sigma(\mathrm{F}, \mathrm{G})\) .

Remark. --- Every neighbourhood of 0 in G in the topology \(\sigma(\mathrm{G},\mathrm{F})\) contains a neighbourhood V defined by a finite number of inequalities of the form \(|\langle x_i,y\rangle |\leqslant 1\) ( \(1\leqslant i\leqslant n\) ), where the \(x_{i}\) are arbitrary points of F. If A is the symmetric convex envelope of the set of the \(x_{i}\) , then V is the polar set \(\mathrm{A}^{\circ}\) of A in G. We can say that the polars in G of finite symmetric sets in F (or of their convex envelopes) form a fundamental system of neighbourhoods of 0 in G for \(\sigma(\mathrm{G},\mathrm{F})\) . If the duality is separating in F, these convex envelopes are compact for \(\sigma(\mathrm{F},\mathrm{G})\) (II, p.~14, cor. 1 of prop. 15), and of finite dimensions. Conversely every compact, convex set of finite dimension C in F (with the \(\sigma(\mathrm{F},\mathrm{G})\) topology) is contained in the convex envelope of a finite subset of F. For, let M be a vector subspace of finite dimension containing C. If \((e_i)_{1\leqslant i\leqslant n}\) is a basis of M, we can suppose that C is contained in the closed parallelotope centre 0 and constructed on the vectors of the basis \(e_i\) (GT, VI, § 1.3); now it is immediate that this parallelotope is the convex envelope of the points \(\sum_{i=1}^{n} \varepsilon_i e_i\) with \(\varepsilon_i = \pm 1\) .

Thus we can say that (if \(\sigma(\mathrm{F},\mathrm{G})\) is Hausdorff) the polars of finite dimensional, convex, compact sets in F (for \(\sigma(\mathrm{F},\mathrm{G})\) or for any Hausdorff locally convex topology finer than \(\sigma(\mathrm{F},\mathrm{G})\) on F) form a fundamental system of neighbourhoods of 0 for \(\sigma(\mathrm{G},\mathrm{F})\) .

COROLLARY 3. --- Let T be the topology of a locally convex space E and let E' be its dual (II, p.~42, def. 1).

\begin{enumerate}
\def\labelenumi{(\roman{enumi})}
\item
  The closed convex sets in \(\mathbf{E}\) are the same for the topology \(\mathcal{T}\) and for the weak topology \(\sigma (\mathrm{E},\mathrm{E}^{\prime})\) .
\item
  For every subset \(\mathbf{M}\) of \(\mathbf{E}\) , the polar set \(\mathbf{M}^{\circ \circ}\) in \(\mathbf{E}\) of the polar set \(\mathbf{M}^{\circ}\) of \(\mathbf{M}\) in \(\mathbf{E}'\) , is the convex closed envelope of \(\mathbf{M} \cup \{0\}\) for the topology \(\mathcal{T}\) .
\end{enumerate}

Clearly, (ii) follows from (i) and th. 1. From the definition of the dual E', it follows from II, p.~43, prop. 3 that the continuous linear forms on E for the topology T are the same as the continuous linear forms for \(\sigma(\mathrm{E}, \mathrm{E}')\) . The closed half-spaces in E are therefore the same for T and for \(\sigma(\mathrm{E}, \mathrm{E}')\) (II, p.~15, prop. 17) and the assertion (i) follows therefore from II, p.~38, cor. 1.

